\documentclass[11pt,aspectratio=169]{ctexbeamer}

\usetheme{default}
\setbeamertemplate{navigation symbols}{}

\usepackage{amsmath}
\usepackage{booktabs}

\definecolor{DeepBlue}{RGB}{25,74,122}
\definecolor{SoftInk}{RGB}{53,59,66}
\definecolor{PaperBg}{RGB}{249,248,245}
\definecolor{PanelBlue}{RGB}{236,243,249}

\setbeamercolor{normal text}{fg=SoftInk,bg=PaperBg}
\setbeamercolor{structure}{fg=DeepBlue}
\setbeamercolor{frametitle}{fg=DeepBlue,bg=PaperBg}
\setbeamercolor{title}{fg=DeepBlue}
\setbeamercolor{block title}{fg=DeepBlue,bg=PanelBlue}
\setbeamercolor{block body}{fg=SoftInk,bg=white}

\newcommand{\Core}[1]{
  \begin{block}{核心结论}
    #1
  \end{block}
}

\newcommand{\LectureMeta}[2]{
  \title{大模型推理系统与实践}
  \subtitle{#1}
  \author{张书豪}
  \institute{华中科技大学计算机学院}
  \date{2027 · #2}
}

\newcommand{\PaperRef}[3]{
  \item \textbf{#1}（#2）\hfill\href{#3}{原文链接}
}

\newcommand{\ResearchSources}[1]{
  \begin{frame}{本讲参考资料}
    \small
    \begin{itemize}
      #1
    \end{itemize}
    \vspace{0.3em}
    \textcolor{SoftInk}{课件中的论文数字均指作者在特定模型、硬件和工作负载下报告的结果，不代表无条件性能。}
  \end{frame}
}

\newcommand{\InternalCaseNote}{
  \small\textcolor{SoftInk}{说明：本页为课程研究案例或教学假设，不代表已发表论文结论；只有在实验包给出模型、硬件、工作负载和原始记录时，数字才可作为结果引用。}
}

% 14 讲细分版的统一参考资料入口。论文数字的解释边界见课程内容基线。

\newcommand{\LectureOneSources}{\ResearchSources{
  \PaperRef{vLLM: Easy, Fast, and Cheap LLM Serving with PagedAttention}{SOSP 2023}{https://doi.org/10.1145/3600006.3613165}
  \PaperRef{SGLang: Efficient Execution of Structured Language Model Programs}{arXiv:2312.07104}{https://arxiv.org/abs/2312.07104}
  \PaperRef{nano-vllm-hust 教学代码}{课程仓库}{https://github.com/intellistream/intro-to-llm-inference-engines}
}}

\newcommand{\LectureTwoSources}{\ResearchSources{
  \PaperRef{vLLM / PagedAttention}{SOSP 2023}{https://doi.org/10.1145/3600006.3613165}
  \PaperRef{Sarathi-Serve}{OSDI 2024}{https://www.usenix.org/conference/osdi24/presentation/agrawal}
  \PaperRef{DistServe}{OSDI 2024}{https://www.usenix.org/conference/osdi24/presentation/zhong-yinmin}
}}

\newcommand{\LectureThreeSources}{\ResearchSources{
  \PaperRef{vLLM / PagedAttention}{SOSP 2023}{https://doi.org/10.1145/3600006.3613165}
  \PaperRef{DistServe}{OSDI 2024}{https://www.usenix.org/conference/osdi24/presentation/zhong-yinmin}
  \PaperRef{Llumnix}{OSDI 2024}{https://www.usenix.org/conference/osdi24/presentation/sun-biao}
}}

\newcommand{\LectureFourSources}{\ResearchSources{
  \PaperRef{Orca}{OSDI 2022}{https://www.usenix.org/conference/osdi22/presentation/yu}
  \PaperRef{Sarathi-Serve}{OSDI 2024}{https://www.usenix.org/conference/osdi24/presentation/agrawal}
  \PaperRef{FastServe}{NSDI 2026}{https://www.usenix.org/conference/nsdi26/presentation/wu-bingyang}
}}

\newcommand{\LectureFiveSources}{\ResearchSources{
  \PaperRef{vLLM / PagedAttention}{SOSP 2023}{https://doi.org/10.1145/3600006.3613165}
  \PaperRef{vAttention}{arXiv:2405.04437}{https://arxiv.org/abs/2405.04437}
  \PaperRef{Rethinking KV Cache Compression}{MLSys 2025}{https://proceedings.mlsys.org/paper_files/paper/2025/hash/26289c647c6828e862e271ca3c490486-Abstract-Conference.html}
  \PaperRef{Mooncake}{FAST 2025}{https://www.usenix.org/conference/fast25/presentation/qin}
}}

\newcommand{\LectureSixSources}{\ResearchSources{
  \PaperRef{DistServe}{OSDI 2024}{https://www.usenix.org/conference/osdi24/presentation/zhong-yinmin}
  \PaperRef{RAGCache}{arXiv:2404.12457}{https://arxiv.org/abs/2404.12457}
  \PaperRef{Mooncake}{FAST 2025}{https://www.usenix.org/conference/fast25/presentation/qin}
  \PaperRef{NEO}{MLSys 2025}{https://proceedings.mlsys.org/paper_files/paper/2025/hash/66a026c0d17040889b50f0dfa650e5e0-Abstract-Conference.html}
}}

\newcommand{\LectureSevenSources}{\ResearchSources{
  \PaperRef{vLLM / PagedAttention}{SOSP 2023}{https://doi.org/10.1145/3600006.3613165}
  \PaperRef{SGLang}{arXiv:2312.07104}{https://arxiv.org/abs/2312.07104}
  \PaperRef{DistServe}{OSDI 2024}{https://www.usenix.org/conference/osdi24/presentation/zhong-yinmin}
  \PaperRef{Mooncake}{FAST 2025}{https://www.usenix.org/conference/fast25/presentation/qin}
}}

\newcommand{\LectureEightSources}{\ResearchSources{
  \PaperRef{FlashAttention}{NeurIPS 2022}{https://arxiv.org/abs/2205.14135}
  \PaperRef{FlexGen}{ICML 2023}{https://proceedings.mlr.press/v202/sheng23a.html}
  \PaperRef{PowerInfer}{SOSP 2024 / arXiv:2312.12456}{https://arxiv.org/abs/2312.12456}
  \PaperRef{FlashInfer}{MLSys 2025}{https://proceedings.mlsys.org/paper_files/paper/2025/hash/dbf02b21d77409a2db30e56866a8ab3a-Abstract-Conference.html}
}}

\newcommand{\LectureNineSources}{\ResearchSources{
  \PaperRef{DEEPSERVE}{USENIX ATC 2025}{https://www.usenix.org/conference/atc25/presentation/hu-junhao}
  \PaperRef{QServe}{MLSys 2025}{https://proceedings.mlsys.org/paper_files/paper/2025/hash/fbe2b2f74a2ece8070d8fb073717bda6-Abstract-Conference.html}
  \PaperRef{vLLM-HUST}{开源维护仓库}{https://github.com/vLLM-HUST/vllm-hust}
}}

\newcommand{\LectureTenSources}{\ResearchSources{
  \PaperRef{vLLM / PagedAttention}{SOSP 2023}{https://doi.org/10.1145/3600006.3613165}
  \PaperRef{Revisiting Disaggregated LLM Serving}{2025 preprint, arXiv:2601.08833}{https://arxiv.org/abs/2601.08833}
  \PaperRef{Rethinking KV Cache Compression}{MLSys 2025}{https://proceedings.mlsys.org/paper_files/paper/2025/hash/26289c647c6828e862e271ca3c490486-Abstract-Conference.html}
}}

\newcommand{\LectureElevenSources}{\ResearchSources{
  \PaperRef{MLPerf Inference}{MLCommons benchmark}{https://mlcommons.org/benchmarks/inference-datacenter/}
  \PaperRef{ACM Artifact Review and Badging}{ACM policy}{https://www.acm.org/publications/policies/artifact-review-and-badging-current}
  \PaperRef{Revisiting Disaggregated LLM Serving}{2025 preprint, arXiv:2601.08833}{https://arxiv.org/abs/2601.08833}
}}

\newcommand{\LectureTwelveSources}{\ResearchSources{
  \PaperRef{nano-vllm-hust}{课程教学代码}{https://github.com/intellistream/intro-to-llm-inference-engines}
  \PaperRef{vLLM-HUST}{HUST 维护的 vLLM fork}{https://github.com/vLLM-HUST/vllm-hust}
  \PaperRef{上游 vLLM}{开源项目}{https://github.com/vllm-project/vllm}
}}

\newcommand{\LectureThirteenSources}{\ResearchSources{
  \PaperRef{课程内容基线与实验材料}{课程仓库}{https://github.com/intellistream/intro-to-llm-inference-engines}
  \PaperRef{DEEPSERVE}{USENIX ATC 2025}{https://www.usenix.org/conference/atc25/presentation/hu-junhao}
  \PaperRef{GhostServe}{MLSys 2026}{https://proceedings.mlsys.org/paper_files/paper/2026/hash/d99e8e80a6c41e148db686918dd7eab3-Abstract-Conference.html}
}}

\newcommand{\LectureFourteenSources}{\ResearchSources{
  \PaperRef{课程完整材料与后续更新}{课程仓库}{https://github.com/intellistream/intro-to-llm-inference-engines}
  \PaperRef{SpecGen}{MLSys 2026}{https://proceedings.mlsys.org/paper_files/paper/2026/hash/66a026c0d17040889b50f0dfa650e5e0-Abstract-Conference.html}
  \PaperRef{GhostServe}{MLSys 2026}{https://proceedings.mlsys.org/paper_files/paper/2026/hash/d99e8e80a6c41e148db686918dd7eab3-Abstract-Conference.html}
}}



\usepackage{tikz}
\usetikzlibrary{arrows.meta,positioning}

\LectureMeta{第 1 讲 课程导论}{研究生课程课件}

\begin{document}

\begin{frame}[plain]
  \titlepage
\end{frame}

\begin{frame}{教学目标}
  \Core{本讲建立课程导论的基本框架，包括研究对象、课程定位与后续 14 讲主线。}
  \begin{itemize}
    \item 区分模型能力、服务能力与系统能力三个分析层次。
    \item 建立“请求 - 状态 - 调度 - 指标”的基本观察框架。
    \item 说明课程讲义、实验与教学代码之间的联动关系。
    \item 给出后续讲次共享的问题地图。
  \end{itemize}
\end{frame}

\begin{frame}{教学样例与课程落点}
  \Core{nano-vllm-hust 在课程中主要作为最小教学样例使用；课程最终目标是支持学生进入 vllm-hust 的优化与开发场景。}
  \begin{itemize}
    \item 前几讲以 nano-vllm-hust 为载体建立请求路径、调度控制点与 KV 状态对象的基本认识。
    \item 该样例用于说明系统主线，而不作为课程长期停留的唯一代码对象。
    \item 后续代码导读、实验扩展与课程项目将逐步过渡到更真实的 vllm-hust 开发语境。
  \end{itemize}
\end{frame}

\begin{frame}{为什么值得独立成课}
  \Core{大模型进入真实在线服务后，研究重点会从“能否生成”转向“能否稳定、经济、可验证地提供服务”。}
  \begin{itemize}
    \item 长上下文、RAG、智能体工作流与多租户负载共同改变了请求形态。
    \item TTFT、P99、goodput、成本与稳定性构成核心服务指标集合。
    \item 系统收益主要取决于全链路组织，而非局部 kernel 或单层算子的单点提速。
  \end{itemize}
\end{frame}

\begin{frame}{从模型问题到系统问题}
  \Core{同一模型在不同运行时、调度策略和状态管理机制下，可能呈现出显著不同的服务质量。}
  \begin{itemize}
    \item 模型层主要关注能力边界、精度与泛化表现。
    \item 系统层主要关注排队、缓存、执行路径、资源治理与回归稳定性。
    \item 因而课程的分析单位是服务系统，而非单一网络层或单篇算法论文。
  \end{itemize}
\end{frame}

\begin{frame}{课程分析对象示意}
  \Core{本课程的分析对象是由请求、状态、控制点与结果指标共同构成的在线推理闭环。}
  \centering
  \begin{tikzpicture}[
    every node/.style={font=\scriptsize, align=center},
    box/.style={draw=DeepBlue, rounded corners, thick, fill=white, minimum width=1.72cm, minimum height=0.9cm},
    line/.style={-Latex, thick, draw=DeepBlue},
    node distance=0.45cm and 0.45cm
  ]
    \node[box] (req) {用户请求\\Prompt / Params};
    \node[box, right=of req] (queue) {请求队列\\Admission};
    \node[box, right=of queue] (sched) {调度器\\Batch / Priority};
    \node[box, right=of sched] (exec) {执行路径\\Prefill / Decode};
    \node[box, right=of exec] (metric) {结果指标\\TTFT / P99 / Goodput};
    \node[box, below=1.0cm of sched] (state) {状态对象\\KV / Memory / Cache};
    \draw[line] (req) -- (queue);
    \draw[line] (queue) -- (sched);
    \draw[line] (sched) -- (exec);
    \draw[line] (exec) -- (metric);
    \draw[line] (sched) -- (state);
    \draw[line] (state) -- (exec);
    \draw[line] (state.west) -| (queue.south);
  \end{tikzpicture}
  \par\vspace{0.5em}
  \raggedright
  \fbox{\parbox{0.9\textwidth}{\small 课程重点是这些对象之间的约束关系，而不是把推理问题缩成一次前向计算。}}
\end{frame}

\begin{frame}{课程边界}
  \Core{本课程关注在线推理服务系统，而不是训练系统或纯算法优化。}
  \begin{itemize}
    \item 分析请求进入、调度、执行与回收的完整路径。
    \item 分析 KV Cache、知识缓存与会话记忆等状态对象。
    \item 分析执行路径、异构平台与实验验证方法。
  \end{itemize}
\end{frame}

\begin{frame}{课程不重点展开的内容}
  \Core{为保持系统问题定义的清晰性，课程将若干相关但非核心议题置于讲授边界之外。}
  \begin{itemize}
    \item 不系统展开预训练、RLHF 与参数高效微调。
    \item 不将提示工程或应用层工作流设计作为主体内容。
    \item 不将单个 kernel 调优等同于整门系统课程。
  \end{itemize}
\end{frame}

\begin{frame}{统一视角：状态化系统}
  \Core{本课程将大模型推理系统表述为“持续携带状态的在线服务系统”。}
  \begin{itemize}
    \item 模型参数之外，还存在 KV Cache、知识缓存、会话记忆与调度元数据。
    \item 状态会跨时间、跨请求与跨设备持续传播。
    \item 调度、缓存与执行路径因此形成连续耦合关系。
  \end{itemize}
\end{frame}

\begin{frame}{状态化系统的含义}
  \Core{状态化系统是指其当前执行决策同时受到当前请求与历史遗留状态的共同约束。}
  \[
    \text{Service Behavior}_t = f(\text{Request}_t, \text{State}_{t-1}, \text{Policy}_t)
  \]
  \begin{itemize}
    \item KV 会影响后续 decode 路径与显存占用。
    \item 会话历史会改变 prompt 长度、状态规模与排队压力。
    \item 运行时统计信息又会反向影响调度与资源治理。
  \end{itemize}
\end{frame}

\begin{frame}{时延分解及其系统含义}
  \Core{首 token 时延与端到端生成时间均由多个系统环节共同决定，而非单一算子性能的直接映射。}
  \[
    \text{TTFT} = T_{\text{queue}} + T_{\text{schedule}} + T_{\text{prefill}} + T_{\text{return}}
  \]
  \[
    T_{\text{end-to-end}} = \text{TTFT} + N_{\text{out}} \cdot T_{\text{decode-step}} + T_{\text{stream}}
  \]
  \begin{itemize}
    \item 若排队或状态驻留构成主导项，单个算子优化通常无法直接改善用户体验。
    \item 这也是后续课程分别讨论工作负载、调度、KV 与验证方法的原因。
  \end{itemize}
\end{frame}

\begin{frame}{课程主线}
  \Core{课程沿着工作负载、调度、状态管理、执行路径和验证方法五条线展开。}
  \begin{enumerate}
    \item 工作负载与评价指标视角
    \item 请求生命周期分解
    \item 调度、KV 与状态管理机制
    \item 异构执行、平台适配与实验方法
    \item 开源系统实践与课程项目
  \end{enumerate}
\end{frame}

\begin{frame}{课程问题地图}
  \Core{整门课都围绕一个总问题展开：如何在动态负载、持续状态和异构约束下构建稳定可验证的推理服务。}
  \begin{itemize}
    \item 工作负载定义问题边界。
    \item 调度决定系统控制点。
    \item 状态组织决定容量与复用空间。
    \item 执行路径决定收益能否兑现。
  \end{itemize}
\end{frame}

\begin{frame}{课程组织方式}
  \Core{课程总评 100 分全部来自 MOD 项目。}
  \begin{itemize}
    \item 问题定义 20 分，实现质量 20 分。
    \item 实验设计 25 分，结果解释 25 分。
    \item 表达与复现 10 分，其中个人问答与贡献 4 分。
    \item 阅读、论文精读、代码导读和最小复现用于准备项目。
  \end{itemize}
\end{frame}

\begin{frame}{四层材料如何配合}
  \Core{整门课按“讲义定问题、tutorial 定路径、实验定证据、代码定控制点”的顺序推进。}
  \begin{itemize}
    \item 讲义与课件首先建立工作负载、调度、状态与执行路径四条主线。
    \item Tutorials 将同一问题先落到 nano-vLLM-HUST 的最小代码对象和日志观察点，再为进入 vllm-hust 做准备。
    \item 课程实验要求学生围绕统一模板提交调用链、表格、日志与可否定解释。
    \item 教学代码同时提供参考实现、student\_exercises、assignment 包与 TA 汇总脚本，以保持讲授、实践与提交的一致性。
  \end{itemize}
\end{frame}

\begin{frame}{课程产出要求}
  \Core{每组开发一个在宿主中可安装、触发和停用的 MOD。}
  \begin{itemize}
    \item 选题：目标问题、扩展点和一条可执行的验收用例。
    \item 联调：代码提交号、打包方式和宿主联调记录。
    \item 交付：插件、测试、对照数据、复现说明和组员分工。
    \item 每位成员在展示时说明贡献并回答技术问题。
  \end{itemize}
\end{frame}

\begin{frame}{课程对象与导论课的统一判断框架}
  \Core{导论部分将研究对象界定为在线推理服务系统，并用对象、控制点、指标和证据构成基本分析框架。}
  \begin{columns}[T,totalwidth=\textwidth]
    \column{0.46\textwidth}
    \begin{block}{课程对象}
      \begin{itemize}
        \item 课堂观察单位不再局限于网络层，而是请求对象、控制点、状态对象与结果指标。
        \item 同一模型在不同 runtime、调度器与 KV 管理策略下，会呈现不同的服务质量。
      \end{itemize}
    \end{block}
    \column{0.5\textwidth}
    \begin{block}{分析框架}
      \begin{enumerate}
        \item 长期存在的系统对象。
        \item 真正改变行为的控制点。
        \item 端到端指标的变化方向。
        \item 支撑解释的日志、表格和代码证据。
      \end{enumerate}
    \end{block}
  \end{columns}
\end{frame}

\begin{frame}{例子：传统部署经验的局限}
  \Core{从分类模型到大模型在线服务，系统主导矛盾已由“单次前向吞吐”转向“跨轮状态组织”。}
  \begin{tabular}{p{0.24\textwidth}p{0.25\textwidth}p{0.35\textwidth}}
    \toprule
    场景 & 主要特征 & 主导系统问题 \\
    \midrule
    传统 CV / NLP 服务 & 一次前向、输入输出短 & 吞吐、批处理、资源利用率 \\
    早期小模型对话 & 有生成过程，但状态较小 & 短任务排队、基础流式返回 \\
    大模型在线服务 & 长上下文、连续 decode、状态驻留 & 调度、状态管理、尾时延与稳定性 \\
    \bottomrule
  \end{tabular}
\end{frame}

\begin{frame}{例子：真实服务中的典型问题}
  \Core{许多服务失效并非源于模型无法生成，而是源于系统对状态与并发的承载能力不足。}
  \begin{tabular}{p{0.25\textwidth}p{0.28\textwidth}p{0.31\textwidth}}
    \toprule
    现象 & 直接表现 & 更可能追查的系统对象 \\
    \midrule
    晚高峰突然超时 & TTFT 和 P99 同时恶化 & 队列、准入、调度策略 \\
    显存未满却无法扩容 & 理论容量未用尽先失去弹性 & KV 块管理、碎片、预留策略 \\
    换一套 runtime 行为大变 & 同一模型服务质量差异明显 & 请求组织、状态路径、执行栈 \\
    \bottomrule
  \end{tabular}
\end{frame}

\begin{frame}{代码例子：从 nano-vllm 到 vllm-hust}
  \Core{nano-vllm-hust 用于建立最小控制路径的理解；相同的分析方法将进一步迁移到 vllm-hust 的优化与开发任务中。}
  \begin{columns}[T,totalwidth=\textwidth]
    \column{0.48\textwidth}
    \begin{block}{入口代码}
      {\ttfamily\small
      \detokenize{from nanovllm import LLM, SamplingParams}\par
      \detokenize{llm = LLM(model_path, enforce_eager=True)}\par
      \detokenize{outputs = llm.generate(prompts, sampling_params)}\par
      }
    \end{block}
    \column{0.48\textwidth}
    \begin{block}{代码观察点与迁移方向}
      \begin{itemize}
        \item \texttt{LLM.generate} 对应接口封装还是更深层 engine 入口。
        \item 请求状态首次被持有的对象位置。
        \item 从 nano-vllm 到 vllm-hust 的方法迁移位置。
      \end{itemize}
    \end{block}
  \end{columns}
\end{frame}

\begin{frame}{代码例子：最小日志的观察价值}
  \Core{导论阶段的插桩应以支撑主路径判断为主要目标，而不以初始阶段覆盖全部细节为目标。}
  \begin{columns}[T,totalwidth=\textwidth]
    \column{0.44\textwidth}
    \begin{block}{最小日志样式}
      {\ttfamily\small
      \detokenize{print(request_id, prompt_len, stage)}\par
      \detokenize{print(batch_size, running_size)}\par
      }
    \end{block}
    \column{0.54\textwidth}
    \begin{block}{观察价值}
      \begin{itemize}
        \item 请求进入、执行与返回这条主路径可以被直接确认。
        \item 阶段日志、状态探针与 benchmark 记录可以在后续实验中逐步增加。
        \item 这也是课程实验组织观察点的基本方式。
      \end{itemize}
    \end{block}
  \end{columns}
\end{frame}

\begin{frame}{常见误读与讨论任务}
  \Core{本页用于区分常见误读，并将讨论收束到对象、控制点与证据三个层面。}
  \begin{columns}[T,totalwidth=\textwidth]
    \column{0.52\textwidth}
    \begin{block}{三种常见误读}
      \begin{itemize}
        \item 输出慢说明模型不够强。
        \item 换更快 kernel 就能解决服务问题。
        \item 设备利用率低就一定是算力浪费。
      \end{itemize}
    \end{block}
    \column{0.44\textwidth}
    \begin{block}{讨论任务}
      \begin{enumerate}
        \item 判断一次“卡顿”更接近排队、状态还是执行路径问题。
        \item 在 \texttt{example.py} 到 \texttt{llm\_engine.py} 之间指出最值得优先阅读的两个函数。
      \end{enumerate}
    \end{block}
  \end{columns}
  \vspace{0.4em}
  \begin{block}{讨论收束问题}
    课程为何先从系统边界切入，而不是直接从单篇机制论文切入。
  \end{block}
\end{frame}


\begin{frame}{课堂讨论}
  \Core{课堂讨论围绕系统边界、用户感知与状态瓶颈三个维度展开。}
  \begin{itemize}
    \item 大模型服务问题为何不能简化为“将模型部署出去”。
    \item 用户感知到的“卡顿”通常对应哪些系统现象。
    \item 哪些状态对象更容易成为瓶颈放大器。
  \end{itemize}
\end{frame}

\begin{frame}{本讲总结}
  \Core{本讲完成了研究对象由“模型”向“服务系统”的转换，并给出了后续讲次共享的分析框架。}
  \begin{itemize}
    \item 研究对象是完整服务链路，而非单一模型。
    \item 研究目标是稳定且可验证的系统收益。
    \item 后续各讲均将回到“状态对象 + 控制点 + 收益链路”的分析框架。
  \end{itemize}
\end{frame}

\LectureOneSources

\end{document}
